\begin{longtable}{ | p{6cm} | p{7.2cm} | p{1.3cm} | p{1cm} | }
\hline\rowcolor{lightgray}
\textbf{名称} & \textbf{描述} &  \textbf{地址} & \textbf{访问} \\ \hline
\endhead
\multicolumn{4}{|l|}{\textbf{配置寄存器}} \\ \hline
\hyperref[regdesc:LPIOMUXGPIONREG]{LP\_IO\_MUX\_GPIO0\_REG} & GPIO0 的 LP IO MUX 配置寄存器 & 0x0000 & R/W \\ \hline
\hyperref[regdesc:LPIOMUXGPIONREG]{LP\_IO\_MUX\_GPIO1\_REG} & GPIO1 的 LP IO MUX 配置寄存器 & 0x0004 & R/W \\ \hline
\hyperref[regdesc:LPIOMUXGPIONREG]{LP\_IO\_MUX\_GPIO2\_REG} & GPIO2 的 LP IO MUX 配置寄存器 & 0x0008 & R/W \\ \hline
\hyperref[regdesc:LPIOMUXGPIONREG]{LP\_IO\_MUX\_GPIO3\_REG} & GPIO3 的 LP IO MUX 配置寄存器 & 0x000C & R/W \\ \hline
\hyperref[regdesc:LPIOMUXGPIONREG]{LP\_IO\_MUX\_GPIO4\_REG} & GPIO4 的 LP IO MUX 配置寄存器 & 0x0010 & R/W \\ \hline
\hyperref[regdesc:LPIOMUXGPIONREG]{LP\_IO\_MUX\_GPIO5\_REG} & GPIO5 的 LP IO MUX 配置寄存器 & 0x0014 & R/W \\ \hline
\hyperref[regdesc:LPIOMUXGPIONREG]{LP\_IO\_MUX\_GPIO6\_REG} & GPIO6 的 LP IO MUX 配置寄存器 & 0x0018 & R/W \\ \hline
\multicolumn{4}{|l|}{\textbf{版本寄存器}} \\ \hline
\hyperref[regdesc:LPIOMUXDATEREG]{LP\_IO\_MUX\_DATE\_REG} & 版本控制寄存器 & 0x01FC & R/W \\ \hline
\end{longtable}
